\section{Methodology}
\label{sec:methodology}

We ask whether the choice of \emph{target formalism} for LLM-generated graph queries interacts with the \emph{adaptation regime} under which the model is used. To isolate this interaction, we hold the reasoning task, the underlying knowledge graph, and the supervision fixed, and vary only two things: the surface language the planner is asked to produce, and whether the planner is prompted or fine-tuned.

\subsection{Problem Formulation}
\label{sec:methodology:formulation}

Given the TITAN knowledge graph $G=(V,E,R)$, typed, bidirectional relations $R$ over nodes $V$, described fully in Section~\ref{sec:benchmark}, and a natural-language question $q$, the task is to produce a program $P$ in a target formalism $F$ such that executing $P$ on $G$ returns an answer set $A \subseteq V$ satisfying $q$. We study two choices of $F$: a compact path DSL introduced by TITAN (Section~\ref{sec:methodology:dsl}), and SPARQL (CITAZIONE), obtained by a validated compilation of the same underlying supervision (Section~\ref{sec:methodology:sparql}). Because both formalisms execute equivalent supervision over equivalent graph representations, any performance gap we observe between $F{=}\text{DSL}$ and $F{=}\text{SPARQL}$ is attributable to the formalism itself and its interaction with the adaptation regime, not to unequal supervision (Figure~\ref{fig:pipeline}).

\begin{figure}[t]
\centering
\begin{tikzpicture}[
  node distance=4mm and 8mm,
  box/.style={draw, rounded corners, align=center, minimum height=7mm,
              inner sep=2.5pt, font=\scriptsize},
  neutralbox/.style={box, fill=titanNeutralLight!50},
  dslbox/.style={box, draw=titanDSL, thick, text=titanDSL, fill=titanDSL!8},
  sparqlbox/.style={box, draw=titanSPARQL, thick, text=titanSPARQL, fill=titanSPARQL!8},
  arr/.style={-{Latex[length=1.8mm]}, thick},
]
\node[neutralbox] (q) {$q$};
\node[neutralbox, right=of q] (planner) {Planner};
\node[dslbox, above right=3mm and 8mm of planner] (dsl) {DSL path};
\node[sparqlbox, below right=3mm and 8mm of planner] (sparql) {SPARQL query};
\coordinate (mid) at ($(dsl)!0.5!(sparql)$);
\node[neutralbox, right=8mm of mid] (ans) {Answer $A$};

\draw[arr] (q) -- (planner);
\draw[arr] (planner) -- (dsl);
\draw[arr] (planner) -- (sparql);
\draw[arr, titanDSL] (dsl) -- (ans);
\draw[arr, titanSPARQL] (sparql) -- (ans);
\end{tikzpicture}
\caption{Same planner input and supervision; only the target formalism differs.}
\label{fig:pipeline}
\end{figure}

\subsection{The TITAN Path DSL}
\label{sec:methodology:dsl}

A DSL program is a step sequence, serialized as \texttt{<PATH>}~step$_1$~\texttt{<SEP>}~\ldots~\texttt{<SEP>}~step$_n$~\texttt{</PATH>}, and each step plays one of five roles. A \emph{relation} step (e.g.\ \texttt{uses\_attack\_pattern}) traverses a typed, bidirectional edge; a \emph{type seed} (\texttt{is\_<type>\_type}) restricts the current frontier to nodes of a given type, or initializes it if none exists yet; \texttt{filter~<keyword>} keeps only the nodes whose associated text satisfies a condition; \texttt{select~<entities>} forks execution into one parallel branch per named entity; and \texttt{exec\_common} / \texttt{exec\_difference} recombine the branches opened by the most recent \texttt{select}, by intersection or symmetric difference respectively. Execution is fully deterministic, which is what lets us score path correctness by direct execution rather than by comparing generated text to a single reference string (Section~\ref{sec:methodology:metrics}).

\subsection{Compiling to SPARQL}
\label{sec:methodology:sparql}

For the comparison to be fair, the SPARQL condition must express the same reasoning content as the DSL.

\paragraph{RDF projection.} $G$ is exported to an RDF graph (N-Triples) that carries over every typed relation edge together with the STIX literal properties (e.g.\ \texttt{name}, \texttt{description}) that SPARQL needs in order to filter and project results. The resulting projection comprises 128{,}026 triples over 27{,}568 distinct subjects and 69 distinct predicates.\footnote{Verified directly against the released \texttt{titan\_graph.nt}; the predicate set includes both the DSL's relation labels and STIX metadata literals the native executor never needs.}

\paragraph{Compilation.} A rule-based compiler maps every gold DSL path to a canonical SPARQL query, mirroring the DSL executor construct-by-construct: entity seeding becomes a bound-subject triple pattern, a relation step a predicate triple pattern, a type seed a type-restricting pattern, \texttt{filter} a \texttt{FILTER} clause, \texttt{select} a set of per-entity variable bindings, and \texttt{exec\_common}/\texttt{exec\_difference} a shared-variable join or a \texttt{MINUS}/\texttt{UNION} combination. These compiled queries supply the few-shot exemplars, the worked examples, and the fine-tuning targets for the SPARQL condition, so the DSL and SPARQL models are exposed to exactly the same underlying reasoning instances.

\paragraph{Validation.} We validate the compiler by executing, on a stratified sample of gold-executable rows, both the native DSL path and its compiled SPARQL counterpart, and comparing the two answer sets independently. The two executions agree on 100.00\% of two independent validation samples (1{,}500 and 4{,}000 rows). This is the basis for treating any DSL/SPARQL gap reported later as a property of the formalism, not an artifact of the comparison.

\subsection{Graph Execution}
\label{sec:methodology:executor}

DSL programs execute deterministically over $G$'s typed adjacency structure; SPARQL queries execute against the RDF projection through a compiled query engine.\footnote{We use \texttt{pyoxigraph} for the throughput it offers at our test-set scale.} Both paths converge on the same output type (a node/entity answer set) which is what every metric in Section~\ref{sec:methodology:metrics} ultimately scores.

\subsection{Prompted and Fine-Tuned Regimes}
\label{sec:methodology:regimes}

We cross three factors: \textbf{formalism} (DSL vs.\ SPARQL), \textbf{adaptation regime} (prompted vs.\ fine-tuned), and \textbf{reasoning style} (NoCoT vs.\ CoT; CITAZIONE), summarized in Table~\ref{tab:design}. The full $2{\times}2{\times}2$ grid is completed on a small set of backbones chosen specifically to isolate the formalism$\times$regime interaction.

\begin{table}[t]
\centering
\small
\begin{tabular}{ll}
\toprule
\textbf{Factor} & \textbf{Levels} \\
\midrule
Formalism & DSL, SPARQL \\
Regime & Prompted, Fine-tuned \\
Reasoning style & NoCoT, CoT \\
\bottomrule
\end{tabular}
\caption{Experimental design: three crossed factors.}
\label{tab:design}
\end{table}

\paragraph{Prompted.} The planner receives a schema description of the relation/predicate vocabulary and node types, derived programmatically from $G$ so that the DSL and SPARQL versions carry identical information content, together with 12 few-shot exemplars stratified by question category and 3 fixed worked examples. The NoCoT and CoT variants of a given prompt are identical except for a single closing instruction sentence, respond with only the program, or reason step by step first, which we ablate in isolation in Appendix~\ref{sec:appendix}.

\paragraph{Fine-tuned.} We fine-tune with LoRA (rank 16, $\alpha{=}16$, dropout 0), learning rate $3\times10^{-4}$ on a linear schedule, effective batch size 16, for up to 8 epochs with early stopping on validation loss.

\paragraph{Backbones.} We evaluate seven instruction-tuned LLMs spanning roughly two orders of magnitude in parameter count: Phi-3.5-mini (3.8B), Qwen2.5-7B, Qwen2.5-7B-Instruct, Qwen2.5-72B, Llama-3.3-70B, GPT-OSS-120B, and DeepSeek-R1-Distill-Llama-70B. Phi-3.5-mini and Qwen2.5-7B are fine-tuned under both formalisms and both reasoning styles; the remaining backbones are evaluated only in the prompted regime. Full per-model configuration is given in Appendix~\ref{sec:appendix}.

\subsection{Evaluation Metrics}
\label{sec:methodology:metrics}

\paragraph{Exact match (EM).} Path EM is the exact match between the predicted and reference DSL step sequences, after whitespace normalization; Query EM is the same comparison between the generated SPARQL text and the canonical query compiled from the reference path (Section~\ref{sec:methodology:sparql}).

\paragraph{Entity-level F1.} We execute the predicted program from the reference starting entities (oracle entity linking) and compare the resulting answer set to the one produced by the reference program, which isolates program quality from entity-resolution noise. F1 is undefined when the reference program itself yields an empty answer set, so we compute it, uniformly across every model and condition, only over gold-executable rows.
